\textbf{Phần 1. Trắc nghiệm lựa chọn}
\setcounterpageref{ex}{0}

\begin{ex} %[2D14V75]
Số gồm $8$ trăm, $5$ chục và $2$ đơn vị được viết là
\choice
{$825$}
{\True $852$}
{$582$}
{$258$}
\loigiai{
Viết là $852$.
\textbf{Đáp án của câu hỏi là B}
}
\end{ex}

\begin{ex} %[2D14V75]
Kết quả của biểu thức $345 + 255 - 180$ là
\choice
{$410$}
{$430$}
{$400$}
{\True $420$}
\loigiai{
$345 + 255 = 600$; $600 - 180 = 420$.
\textbf{Đáp án của câu hỏi là D}
}
\end{ex}

\begin{ex} %[2D14V75]
Tìm $x$ biết: $x \times 5 = 45$.
\choice
{\True $9$}
{$8$}
{$7$}
{$10$}
\loigiai{
$x = 45 : 5 = 9$.
\textbf{Đáp án của câu hỏi là A}
}
\end{ex}

\begin{ex} %[2D14V75]
Một mảnh vườn hình chữ nhật có chiều dài $24\text{ m}$, chiều rộng $16\text{ m}$. Chu vi mảnh vườn đó là
\choice
{$40\text{ m}$}
{$70\text{ m}$}
{\True $80\text{ m}$}
{$90\text{ m}$}
\loigiai{
Chu vi: $(24 + 16) \times 2 = 80$ m.
\textbf{Đáp án của câu hỏi là C}
}
\end{ex}

\begin{ex} %[2D14V75]
Số tròn trăm lớn nhất có ba chữ số đem trừ đi số chẵn bé nhất có ba chữ số thì được
\choice
{$798$}
{\True $800$}
{$899$}
{$802$}
\loigiai{
Số tròn trăm lớn nhất có 3 chữ số: 900. Số chẵn bé nhất có 3 chữ số: 100. Hiệu: $900 - 100 = 800$.
\textbf{Đáp án của câu hỏi là B}
}
\end{ex}

\begin{ex} %[2D14V75]
Một người có 1 tờ 1 000 đồng, mua 2 quyển vở giá 400 đồng một quyển. Người đó còn lại bao nhiêu đồng?
\choice
{$100\text{ đồng}$}
{$300\text{ đồng}$}
{$400\text{ đồng}$}
{\True $200\text{ đồng}$}
\loigiai{
Mua hết: $400 \times 2 = 800$ đồng. Còn lại: $1 000 - 800 = 200$ đồng.
\textbf{Đáp án của câu hỏi là D}
}
\end{ex}

\begin{ex} %[2D14V75]
Một sợi dây đồng dài $1\text{ m}$. Người ta cắt đi $35\text{ cm}$. Chiều dài đoạn dây còn lại là
\choice
{\True $65\text{ cm}$}
{$75\text{ cm}$}
{$55\text{ cm}$}
{$65\text{ dm}$}
\loigiai{
$1\text{ m} = 100\text{ cm}$. Còn lại: $100 - 35 = 65$ cm.
\textbf{Đáp án của câu hỏi là A}
}
\end{ex}

\begin{ex} %[2D14V75]
Điền dấu thích hợp: $5 \times 8 \dots 100 - 60$.
\choice
{$>$}
{$<$}
{\True $=$}
{Không so sánh được}
\loigiai{
$5 \times 8 = 40$; $100 - 60 = 40$. Hai vế bằng nhau.
\textbf{Đáp án của câu hỏi là C}
}
\end{ex}

\begin{ex} %[2D14V75]
Đường gấp khúc gồm 4 đoạn thẳng bằng nhau, mỗi đoạn dài 25 cm. Độ dài đường gấp khúc là bao nhiêu mét?
\choice
{$2\text{ m}$}
{\True $1\text{ m}$}
{$10\text{ m}$}
{$100\text{ m}$}
\loigiai{
$25 \times 4 = 100\text{ cm} = 1\text{ m}$.
\textbf{Đáp án của câu hỏi là B}
}
\end{ex}

\begin{ex} %[2D14V75]
Khối 2 có $280$ học sinh, khối 3 có nhiều hơn khối 2 là $45$ học sinh. Cả hai khối có tất cả bao nhiêu học sinh?
\choice
{$595$ học sinh}
{$615$ học sinh}
{$325$ học sinh}
{\True $605$ học sinh}
\loigiai{
Khối 3: $280 + 45 = 325$. Cả hai khối: $280 + 325 = 605$ học sinh.
\textbf{Đáp án của câu hỏi là D}
}
\end{ex}

\begin{ex} %[2D14V75]
Tìm một số biết rằng nếu lấy số đó chia cho 2 rồi cộng với 15 thì được 24.
\choice
{\True $18$}
{$16$}
{$20$}
{$22$}
\loigiai{
Số đó chia 2 được: $24 - 15 = 9$. Số đó là: $9 \times 2 = 18$.
\textbf{Đáp án của câu hỏi là A}
}
\end{ex}

\begin{ex} %[2D14V75]
Trong các khẳng định sau, khẳng định nào ĐÚNG?
\choice
{1 m = 10 cm}
{1 ngày = 12 giờ}
{\True 1 km = 1 000 m}
{100 đồng > 200 đồng}
\loigiai{
$1\text{ km} = 1 000\text{ m}$ là khẳng định chính xác.
\textbf{Đáp án của câu hỏi là C}
}
\end{ex}

\textbf{Phần 2. Trắc nghiệm đúng sai}
\setcounterpageref{ex}{0}

\begin{ex} %[2D14V75]
Xét các bài toán tổng hợp toàn diện chương trình Toán 2:
\choiceTF
{\True Số chẵn liền sau của 998 là 1 000.}
{Hình tam giác có 4 cạnh.}
{\True 500 ml + 500 ml = 1 l (hoặc 500 g + 500 g = 1 kg).}
{\True Phép nhân 2 x 9 có kết quả bằng phép nhân 9 x 2.}
\loigiai{
\begin{itemize}
    \item Ý a) Đúng: Đúng.
    \item Ý b) Sai: Hình tam giác có 3 cạnh. Sai.
    \item Ý c) Đúng: Đúng.
    \item Ý d) Đúng: Tính chất giao hoán. Đúng.
\end{itemize}
\begin{itemize}
    \item a) \textbf{Đúng}.
    \item b) \textbf{Sai}.
    \item c) \textbf{Đúng}.
    \item d) \textbf{Đúng}.
\end{itemize}
}
\end{ex}

\begin{ex} %[2D14V75]
Một kho hàng có 1 000 thùng sữa. Ngày đầu xuất đi 350 thùng, ngày thứ hai xuất đi 450 thùng:
\choiceTF
{\True Cả hai ngày xuất đi 800 thùng sữa.}
{\True Trong kho còn lại 200 thùng sữa.}
{Số sữa còn lại nhiều hơn số sữa đã xuất.}
{\True Nếu mỗi thùng xếp được 5 hộp sữa thì 200 thùng còn lại có 1 000 hộp sữa.}
\loigiai{
\begin{itemize}
    \item Ý a) Đúng: 350 + 450 = 800 thùng. Đúng.
    \item Ý b) Đúng: 1 000 - 800 = 200 thùng. Đúng.
    \item Ý c) Sai: 200 < 800 nên ít hơn. Sai.
    \item Ý d) Đúng: 200 x 5 = 1 000. Đúng.
\end{itemize}
\begin{itemize}
    \item a) \textbf{Đúng}.
    \item b) \textbf{Đúng}.
    \item c) \textbf{Sai}.
    \item d) \textbf{Đúng}.
\end{itemize}
}
\end{ex}

\textbf{Phần 3. Trả lời ngắn}
\setcounterpageref{ex}{0}

\begin{ex} %[2D14V75]
Tính giá trị của biểu thức: $425 + 375 - 200$.
\shortans{600}
\loigiai{
\textbf{Đáp số:} $600$\\
$800 - 200 = 600$.
}
\end{ex}

\begin{ex} %[2D14V75]
Tìm $x$ biết: $x : 5 + 12 = 18$.
\shortans{30}
\loigiai{
\textbf{Đáp số:} $30$\\
$x : 5 = 6 \Rightarrow x = 30$.
}
\end{ex}

\begin{ex} %[2D14V75]
Một hình tam giác có ba cạnh bằng nhau, chu vi bằng 24 cm. Mỗi cạnh dài bao nhiêu xăng-ti-mét?
\shortans{8}
\loigiai{
\textbf{Đáp số:} $8$\\
$24 : 3 = 8$ cm.
}
\end{ex}

\begin{ex} %[2D14V75]
An có 1 000 đồng, An mua một cây kem hết 800 đồng. An còn lại bao nhiêu đồng?
\shortans{200}
\loigiai{
\textbf{Đáp số:} $200$\\
$1 000 - 800 = 200$ đồng.
}
\end{ex}

\begin{ex} %[2D14V75]
Một xe khách chở 45 người. Đến trạm dừng có 12 người xuống và 7 người lên. Trên xe còn lại bao nhiêu người?
\shortans{40}
\loigiai{
\textbf{Đáp số:} $40$\\
$45 - 12 + 7 = 40$ người.
}
\end{ex}

\begin{ex} %[2D14V75]
Tính chu vi hình chữ nhật có chiều dài 30 cm và chiều rộng 20 cm.
\shortans{100}
\loigiai{
\textbf{Đáp số:} $100$\\
$(30 + 20) \times 2 = 100$ cm.
}
\end{ex}
